\documentclass{itmo}

\setlist[itemize]{label=---}
\setlist[enumerate]{label=\arabic*)}

\usepackage{dcolumn}
\newcolumntype{d}[1]{D{.}{.}{#1}}

\begin{document}

\makereporttitle{Программной инженерии и компьютерной техники (ФПИ и КТ)}
	{лабораторной работе №1}
	{Структуры и алгоритмы в базах данных и распределенных системах}
	{Реализация алгоритмов хэширования (perfect, extendible, min)}
	{}
	{Постнов~С.~А./P4135}
	{Платонов~А.~В./доцент, Портнов~П.~В./ассистент}

\maketableofcontents

\chapter{Теоретическая часть}

\section{Perfect hashing}

Задача: хранить пары ключ--значение так, чтобы операции \texttt{set/get} выполнялись за $O(1)$ (в среднем) и не возникало коллизий для фиксированного набора ключей.

\textit{Perfect hash} для множества ключей $S$~---~это хэш--функция $h$, которая отображает ключи из $S$ в индексы таблицы без коллизий.
В классическом варианте (FKS) используется двухуровневая схема: первичная таблица распределяет ключи по корзинам, а для каждой корзины строится вторичная таблица, в которой подбираются параметры так, чтобы коллизий не было.

В данной работе реализована схожая идея:
\begin{itemize}
	\item первичная таблица имеет размер, равный ближайшему простому $\ge n$;
	\item в ячейке хранится либо одна пара \texttt{key/value}, либо указатель на вторичную таблицу;
	\item при коллизии создаётся/перестраивается вторичная таблица: увеличивается её размер и перебираются различные хэш--функции до получения раскладки без коллизий.
\end{itemize}

\chapter{Практическая часть}

\section{Реализация perfect hash таблицы}

Публичный интерфейс реализован в файлах \texttt{inc/perfect.h} и \texttt{src/perfect.c}:
\begin{itemize}
	\item \texttt{create\_perfect\_table} создаёт таблицу нужного размера;
	\item \texttt{perfect\_set} добавляет/обновляет пару ключ--значение;
	\item \texttt{perfect\_get} выполняет поиск значения по ключу;
	\item \texttt{free\_perfect\_table} освобождает память.
\end{itemize}

Структура \texttt{perfect\_hash\_table\_t} содержит массив указателей на \texttt{ceil\_value\_t}. Каждая ячейка либо хранит одну пару \texttt{key/value}, либо указывает на вторичную таблицу (вложенный \texttt{perfect\_hash\_table\_t}).

Для хэширования строк используются функции из \texttt{inc/hash.h}. При перестроении вторичных таблиц перебирается набор хэш--функций до построения раскладки без коллизий.

\chapter{Исследовательская часть}

\section{Методика замеров}

Для замеров использована программа \texttt{src/benchmark.c}. Для каждого датасета из \texttt{data/bench/manifest.txt}:
\begin{enumerate}
	\item создаётся пустая perfect hash таблица;
	\item измеряется время вставки всех $n$ ключей (\texttt{perfect\_set});
	\item измеряется время чтения всех $n$ ключей (\texttt{perfect\_get}).
\end{enumerate}
Каждый прогон усредняется по \texttt{BENCH\_ITERS} (по умолчанию 5).

\section{Результаты}

Таблица~\ref{tab:perfect-bench} содержит усреднённые значения задержек и пропускной способности.
\begin{table}[h]
	\centering
	\caption{Результаты бенчмарка perfect hash (среднее по итерациям)}
	\label{tab:perfect-bench}
	\begin{tabularx}{\textwidth}{|X|X|X|X|X|}
		\toprule
		$N$ & вставка, нс/оп & вставка, млн оп/с & поиск, нс/оп & поиск, млн оп/с \\
		\midrule
		50 & 344.0 & 2.98 & 80.0 & 12.50 \\
		250 & 528.0 & 1.90 & 112.0 & 8.95 \\
		1250 & 226.4 & 4.42 & 93.0 & 10.76 \\
		7500 & 314.9 & 3.24 & 82.8 & 12.19 \\
		20000 & 404.5 & 2.51 & 66.2 & 15.22 \\
		50000 & 217.0 & 4.61 & 52.4 & 19.09 \\
		100000 & 127.7 & 7.83 & 43.2 & 23.15 \\
		250000 & 225.7 & 4.43 & 58.3 & 17.14 \\
		500000 & 141.3 & 7.08 & 52.2 & 19.15 \\
		750000 & 138.7 & 7.21 & 52.0 & 19.24 \\
		1000000 & 197.4 & 5.07 & 57.9 & 17.29 \\
		2000000 & 164.9 & 6.07 & 57.2 & 17.51 \\
		5000000 & 194.1 & 5.15 & 81.2 & 12.32 \\
		\bottomrule
	\end{tabularx}
\end{table}

Графики результатов замеров представлены на рисунках~\ref{img:bench_avg_latency} и~\ref{img:bench_throughput}:

\includeimage
	{bench_avg_latency}
	{f}
	{H}
	{\textwidth}
	{Средняя задержка операций вставки/поиска для perfect hash}

\includeimage
	{bench_throughput}
	{f}
	{H}
	{\textwidth}
	{Пропускная способность операций вставки/поиска для perfect hash}

\chapter{Вывод}

Реализована perfect hash таблица для строковых ключей с поддержкой операций вставки и поиска.
Проведены замеры производительности на наборах данных разного размера; результаты приведены в таблице и на графиках.

\end{document}
